% !TeX program = pdflatex
% =============================================================================
% Aufgabenblatt: Ohmsches Gesetz (Strom, Spannung, Widerstand)
% UZH Praktikum I -- Sara Romer Urban, Kantonsschule im Lee, HS2026
% Klasse: 1f (16.09.2026), aus PLAN_Dokumentengenerierung.md.
%
% ENTWURF, NOCH NICHT INS REPO KOPIERT (auf Wunsch von Loris).
%
% Gehoert zusammen mit arbeitsblatt3-u-i-r.tex im selben Ordner: dort steht
% die gesamte Theorie (Einheiten, Ohmsches Gesetz, Umformen, Vorsaetze),
% hier ausschliesslich die Uebungsaufgaben dazu -- Aufteilung nach Loris'
% eigenem Kriterium (Arbeitsblatt = Plenum, Aufgabenblatt = Aufgaben), nicht
% nach dem generischen Werkzeug-Kriterium aus institutionen/
% UZH-Praktikum-I/CLAUDE.md. Deshalb bewusst OHNE einleitende Erklaerbox,
% beginnt direkt mit der ersten Aufgabe.
%
% Reihenfolge: drei einfache Anwendungen des Ohmschen Gesetzes (je einmal
% nach I, R, U aufgeloest), danach reine Vorsatz-Umrechnung (mA/A, kOhm/Ohm,
% mV/V), zum Schluss eine Aufgabe, die Ohmsches Gesetz und Vorsatz-
% Umrechnung kombiniert. Kein Bezug auf ein U-I-Diagramm.
%
% Numerierung (Aufgabenblatt-Zaehlung nach institutionen/UZH-Praktikum-I/
% CLAUDE.md) noch offen -- Dateiname bewusst ohne Zahl, bis die naechste
% freie Nummer bekannt ist (siehe dortiger Abschnitt "Nummerierung und
% Dateibenennung").
%
% -----------------------------------------------------------------------------
% Korrektur (16.09.2026, Abgleich mit den zwei anderen Blaettern zum
% selben Thema, gruppenpuzzle3-u-i-r.tex und arbeitsblatt3-u-i-r.tex):
% Inhaltlich/terminologisch passt dieses Aufgabenblatt bereits zum
% Vorwissen aus den beiden anderen Dokumenten (Verbraucher-Begriff,
% Widerstand/Spannung/Strom-Notation, nur Milli/Kilo als Vorsaetze, keine
% Serie-/Parallelschaltung-Rechnung, da Ersatzwiderstand dort bewusst noch
% nicht eingefuehrt wird) -- keine Aenderung noetig. Behoben wurde aber
% derselbe technische Fehler, der in beiden anderen Dokumenten bereits
% aufgefallen und dort per \setcounter{question}{...} korrigiert war: Die
% drei separaten \begin{questions}-Bloecke (Übung 1-3) liessen den
% exam-Klassen-Fragezaehler bei jedem Block neu bei 1 beginnen (Warnung
% "Label question@1 multiply defined" im Log), statt fortlaufend 1-5 zu
% zaehlen. Jetzt mit \setcounter{question}{3} bzw. {4} vor Übung 2/3 auf
% durchgehende Nummerierung 1-5 korrigiert, analog zu den anderen beiden
% Dokumenten.
% =============================================================================
\documentclass[addpoints,11pt,a4paper]{exam}

\usepackage[T1]{fontenc}
\usepackage[utf8]{inputenc}
\usepackage[ngerman]{babel}
\usepackage{lmodern}
\usepackage[a4paper,margin=2.2cm,headheight=14pt]{geometry}
\usepackage{amsmath}
\usepackage{siunitx}
% Hausstil: zusammengesetzte Einheiten immer als echter Bruch (\frac), nicht
% als Inline-Schraegstrich oder \tfrac -- gilt global fuer jedes \si/\SI.
\sisetup{per-mode=fraction,fraction-command=\frac}
\usepackage{array,booktabs,tabularx,multirow}
\usepackage{enumitem}
\usepackage{xcolor}
\usepackage{colortbl}
\usepackage[most]{tcolorbox}
\usepackage{lastpage}
\usepackage{graphicx}
\usepackage{hyperref}

\usepackage{setspace}
\usepackage{needspace}
\setlength{\parindent}{0pt}
\setlength{\parskip}{0.6em}
\setstretch{1.08}
\setlist[itemize]{itemsep=0.4em}

% -----------------------------
% Farben (Corporate-Farbschema Physik KFR)
% -----------------------------
\definecolor{titleblue}{HTML}{183A6B}
\definecolor{accentblue}{HTML}{2E6F9E}
\definecolor{darkgray}{HTML}{4A4A4A}

% Links (hyperref): farblich ans Corporate-Farbschema angeglichen.
\hypersetup{
  colorlinks=true,
  urlcolor=accentblue,
  linkcolor=accentblue,
  pdftitle={Aufgabenblatt: Ohmsches Gesetz},
  pdfauthor={Loris Keller}
}

% -----------------------------
% Boxen
% -----------------------------
\tcbset{before skip=10pt, after skip=10pt}
\newtcolorbox{titlebox}{
  enhanced,
  colback=titleblue,
  colframe=titleblue,
  boxrule=0pt,
  arc=3mm,
  left=3mm,right=3mm,top=3mm,bottom=3mm
}

\newlength{\aufgabenindent}
\settowidth{\aufgabenindent}{10.\hskip\labelsep}

\newtcolorbox{taskbox}[1][]{
  enhanced, breakable,
  colback=white, colframe=white, boxrule=0pt,
  left=0mm,right=0mm,top=0mm,bottom=0mm,
  before skip=10pt, after skip=10pt,
  before upper={%
    \def\tmpboxtitle{#1}%
    \ifx\tmpboxtitle\empty\else\noindent\textbf{#1}\par\smallskip\fi
  }
}

\newtcolorbox{answerbox}{
  breakable,
  colback=white,
  colframe=gray!55,
  boxrule=0.5pt,
  arc=1.5mm,
  left=2mm,right=2mm,top=2mm,bottom=2mm
}

% --- Loesungen ein-/ausblenden -----------------------------------------------
%\noprintanswers
\printanswers

\pointformat{}
\renewcommand{\solutiontitle}{\noindent\color{titleblue}\textbf{L\"osung:}\enspace}

\pagestyle{headandfoot}
\cfoot{\textcolor{darkgray}{Seite \thepage\ von \pageref{LastPage}}}

\newcommand{\Kopfzeile}[4]{%
  \begin{titlebox}
    {\color{white}\sffamily\bfseries\Large #4}
  \end{titlebox}
  \vspace{0.5em}
  \noindent\textbf{Klasse:}~#1 \hfill \textbf{Datum:}~#2 \hfill \textbf{Thema:}~#3
    \hfill \textbf{Lehrperson:}~Loris Keller
  \vspace{0.8em}
  \lhead{\textcolor{darkgray}{#3}}
  \rhead{\textcolor{darkgray}{#4}}
}

\newlength{\antwortraumlen}
\newcommand{\Antwortraum}[1]{%
  \pgfmathsetlength{\antwortraumlen}{#1*2.5cm}%
  \begin{answerbox}
    \rule{0pt}{\antwortraumlen}%
  \end{answerbox}%
}

% Werte fuer Kopfzeile zentral setzen
\newcommand{\Klasse}{1f}
\newcommand{\Datum}{16.09.2026}
\newcommand{\Thema}{Ohmsches Gesetz}
\newcommand{\ArbeitsblattTitel}{Aufgabenblatt: Ohmsches Gesetz}

\begin{document}

\Kopfzeile{\Klasse}{\Datum}{\Thema}{\ArbeitsblattTitel}

% =============================================================================
% Übung 1: Das Ohmsche Gesetz anwenden
% =============================================================================
\begin{questions}
\begin{taskbox}[Übung: Das Ohmsche Gesetz anwenden]
\question[2] An einem Widerstand von $R=\SI{100}{\ohm}$ liegt eine Spannung
von $U=\SI{6}{\volt}$ an. Wie gross ist der Strom $I$, der durch ihn
fliesst?
\Antwortraum{2}
\begin{solution}
\[
  I = \frac{U}{R} = \frac{\SI{6}{\volt}}{\SI{100}{\ohm}} = \SI{0.06}{\ampere}.
\]
\end{solution}

\question[2] Durch eine Lampe fliesst ein Strom von $I=\SI{0.5}{\ampere}$,
an ihr liegt eine Spannung von $U=\SI{12}{\volt}$ an. Welchen Widerstand
hat die Lampe?
\Antwortraum{2}
\begin{solution}
\[
  R = \frac{U}{I} = \frac{\SI{12}{\volt}}{\SI{0.5}{\ampere}} = \SI{24}{\ohm}.
\]
\end{solution}

\question[2] Ein Widerstand von $R=\SI{470}{\ohm}$ wird von einem Strom
von $I=\SI{0.02}{\ampere}$ durchflossen. Welche Spannung fällt über ihm
ab?
\Antwortraum{2}
\begin{solution}
\[
  U = R\cdot I = \SI{470}{\ohm}\cdot\SI{0.02}{\ampere} = \SI{9.4}{\volt}.
\]
\end{solution}
\end{taskbox}
\end{questions}

% =============================================================================
% Übung 2: Vorsätze umrechnen
% =============================================================================
\begin{questions}
\setcounter{question}{3}
\begin{taskbox}[Übung: Vorsätze umrechnen]
\question[4] Rechnen Sie um.
\begin{parts}
  \part[1] $\SI{25}{\milli\ampere} = \underline{\hspace{2.5cm}}\ \si{\ampere}$
  \begin{solution} $\SI{0.025}{\ampere}$ \end{solution}
  \part[1] $\SI{0.008}{\ampere} = \underline{\hspace{2.5cm}}\ \si{\milli\ampere}$
  \begin{solution} $\SI{8}{\milli\ampere}$ \end{solution}
  \part[1] $\SI{3.3}{\kilo\ohm} = \underline{\hspace{2.5cm}}\ \si{\ohm}$
  \begin{solution} $\SI{3300}{\ohm}$ \end{solution}
  \part[1] $\SI{150}{\milli\volt} = \underline{\hspace{2.5cm}}\ \si{\volt}$
  \begin{solution} $\SI{0.15}{\volt}$ \end{solution}
\end{parts}
\end{taskbox}
\end{questions}

% =============================================================================
% Übung 3: Ohmsches Gesetz und Vorsätze kombiniert
% =============================================================================
\begin{questions}
\setcounter{question}{4}
\begin{taskbox}[Übung: Ohmsches Gesetz und Vorsätze kombiniert]
\question[4] In einer Schaltung liegt ein Widerstand von $R=\SI{470}{\ohm}$
in Serie mit einem Verbraucher. Durch den Widerstand fliesst dabei ein
Strom von $I=\SI{15}{\milli\ampere}$.
\begin{parts}
  \part[3] Berechnen Sie die Spannung, die über dem Widerstand abfällt.
  Geben Sie das Ergebnis sowohl in Volt als auch in Millivolt an.
  \Antwortraum{3}
  \begin{solution}
  \[
    U = R\cdot I = \SI{470}{\ohm}\cdot\SI{0.015}{\ampere} = \SI{7.05}{\volt} = \SI{7050}{\milli\volt}.
  \]
  \end{solution}
  \part[1] Eine 9-V-Batterie liegt an diesem Widerstand an (statt der
  Spannung aus Teilaufgabe a). Ist der Strom durch den Widerstand dann
  grösser oder kleiner als $\SI{15}{\milli\ampere}$? Begründen Sie kurz,
  ohne zu rechnen.
  \Antwortraum{1}
  \begin{solution}
  Grösser: $\SI{9}{\volt}$ ist mehr als die $\SI{7.05}{\volt}$ aus
  Teilaufgabe a. Da $R$ gleich bleibt, muss nach $I=U/R$ auch $I$
  zunehmen, wenn $U$ zunimmt.
  \end{solution}
\end{parts}
\end{taskbox}
\end{questions}

\end{document}
